% =====================================================================
%  Chapter 5 — The Chernoff Bound
%  File: chapters/05_chernoff.tex
% =====================================================================
\providecommand{\E}{\mathbb{E}}
\providecommand{\Prob}{\mathbb{P}}
\providecommand{\Var}{\operatorname{Var}}
\providecommand{\whp}{\textit{whp}}

\chapter{The Chernoff Bound}
\label{chap:chernoff}

Chebyshev bounds a deviation by the variance, and decays only polynomially. For
a sum of independent bounded variables the tail is far smaller --- exponential,
on the Gaussian scale. The Chernoff bound captures this, and it underlies most
of the concentration arguments that follow.

\begin{theorem}{Chernoff bound}{chernoff-pm}
Let $S_n=X_1+\cdots+X_n$ with the $X_i$ independent and uniform on $\{-1,+1\}$.
Then for every $\lambda>0$,
\[
  \Prob\bigl(S_n\ge\lambda\sqrt n\bigr)\le e^{-\lambda^2/2},
\]
and the same bound holds for $\Prob(S_n\le-\lambda\sqrt n)$.
\end{theorem}

\begin{proof}
For $t\ge0$, independence factorizes the moment generating function:
\[
  \E\bigl[e^{tS_n}\bigr]=\prod_{i=1}^n\E\bigl[e^{tX_i}\bigr]
  =\Bigl(\frac{e^{t}+e^{-t}}2\Bigr)^n=(\cosh t)^n\le e^{nt^2/2},
\]
where $\cosh t=\sum_{k\ge0}t^{2k}/(2k)!\le\sum_{k\ge0}t^{2k}/(2^k k!)=e^{t^2/2}$.
Markov applied to $e^{tS_n}$ gives
$\Prob(S_n\ge\lambda\sqrt n)\le e^{-t\lambda\sqrt n}\E[e^{tS_n}]\le
e^{nt^2/2-t\lambda\sqrt n}$; the choice $t=\lambda/\sqrt n$ minimizes the exponent
and yields $e^{-\lambda^2/2}$.
\end{proof}

The same moment-generating-function method covers any independent bounded
summands.

\begin{theorem}{Chernoff for bounded variables}{chernoff-bounded}
Let $X_1,\dots,X_n$ be independent with $X_i\in[-1,1]$ and $\E[X_i]=0$. Then
$S_n=\sum_i X_i$ satisfies $\Prob(S_n\ge\lambda\sqrt n)\le e^{-\lambda^2/2}$, and
likewise for the lower tail. In particular, a sum $X$ of independent Bernoullis
with mean $\mu$ obeys $\Prob(|X-\mu|\ge\lambda\sqrt n)\le 2e^{-\lambda^2/2}$.
\end{theorem}

\begin{insight}
Two comparisons place the bound. Against Chebyshev's $1/\lambda^2$, the
exponential $e^{-\lambda^2/2}$ is what lets a union bound survive over
exponentially many events. Against the Gaussian it is essentially tight: the
same computation gives $\Prob(Z\ge\lambda)\le e^{-\lambda^2/2}$ for
$Z\sim\mathcal N(0,1)$, off the true tail only by a factor $\sqrt{2\pi}\lambda$.
The cost is independence --- Chernoff needs the full product structure to factor
the moment generating function, where Chebyshev asks only for pairwise
independence.
\end{insight}

% =====================================================================
\section{Discrepancy}
% =====================================================================

Given a set system $S_1,\dots,S_m$ on a ground set $[n]$, a two-colouring
$\chi\colon[n]\to\{-1,+1\}$ has \emph{discrepancy}
$\operatorname{disc}(\chi)=\max_j\bigl|\sum_{i\in S_j}\chi(i)\bigr|$; the
discrepancy of the system is the minimum over $\chi$. We ask how balanced a
colouring can be forced to exist.

\begin{theorem}{Random colouring discrepancy}{discrepancy}
Any system of $m$ sets on $[n]$ has a two-colouring of discrepancy at most
$\sqrt{2n\ln(2m)}$.
\end{theorem}

\begin{proof}
Colour each point independently and uniformly. For a fixed set $S_j$ the sum
$\sum_{i\in S_j}\chi(i)$ is a $\pm1$ sum of $|S_j|\le n$ terms, so by
Theorem~\ref{thm:chernoff-pm},
$\Prob\bigl(|\sum_{i\in S_j}\chi(i)|\ge\lambda\sqrt n\bigr)\le2e^{-\lambda^2/2}$.
Taking $\lambda=\sqrt{2\ln(2m)}$ makes each bound $1/m$, and a union bound over
the $m$ sets leaves positive probability that every set is balanced within
$\lambda\sqrt n=\sqrt{2n\ln(2m)}$.
\end{proof}

For the balanced case $m=n$ this gives $O(\sqrt{n\log n})$. Remarkably, the
logarithm is removable.

\begin{theorem}{Six standard deviations suffice; Spencer 1985}{spencer}
Any system of $n$ sets on $n$ points has a two-colouring of discrepancy at most
$6\sqrt n$.
\end{theorem}

\begin{insight}
The $\sqrt{\log n}$ gap between random and optimal is real but avoidable. The
union bound treats the $n$ sets as independent adversaries and pays
$\sqrt{\log n}$ to control them together; Spencer's theorem removes the tax, but
not through any single random colouring. Its proof colours a constant fraction
of the points at a time with discrepancy $O(\sqrt n)$ and iterates --- an early
use of the pigeonhole-and-entropy method now standard in discrepancy theory.
Whether the same holds when each point lies in few sets is the Beck--Fiala and
Koml\'os problem, still open.
\end{insight}

% =====================================================================
\section{Nearly orthogonal vectors}
% =====================================================================

How many unit vectors in $\mathbb R^n$ can be pairwise \emph{nearly} orthogonal?
In the exact case at most $n$ vectors can be orthonormal, but relaxing to small
inner products changes the answer dramatically.

\begin{theorem}{Exponentially many nearly orthogonal vectors}{equiangular}
There is a constant $c>0$ such that for every $n$ one can find
$e^{cn}$ vectors in $\{-1,+1\}^n$ whose pairwise inner products all satisfy
$|\langle u,v\rangle|\le C\sqrt{n\log(1/\!\cos)}$; normalizing by $1/\sqrt n$
yields $e^{cn}$ unit vectors with all pairwise inner products $O(\sqrt{(\log)/n})$.
\end{theorem}

\begin{proof}[Proof idea]
Take $N=e^{cn}$ independent uniform vectors in $\{-1,+1\}^n$. For a fixed pair,
$\langle u,v\rangle=\sum_i u_iv_i$ is a $\pm1$ sum of $n$ independent terms, so
by Theorem~\ref{thm:chernoff-pm} it exceeds $\lambda\sqrt n$ in absolute value
with probability at most $2e^{-\lambda^2/2}$. Union-bounding over the
$\binom N2<e^{2cn}$ pairs with $\lambda=\Theta(\sqrt n)$ leaves positive
probability that every pair is nearly orthogonal; choosing $c$ small keeps the
union bound below $1$.
\end{proof}

\begin{insight}
Exact orthogonality caps you at $n$ vectors; near-orthogonality is a much softer
condition that random sign vectors satisfy exponentially often, because two of
them have a normalized inner product that is a mean of $n$ signs and hence
small. This is concentration of measure in one line, and the same fact drives
Johnson--Lindenstrauss dimension reduction.
\end{insight}

% =====================================================================
\section{A counterexample to the Haj\'os conjecture}
% =====================================================================

Haj\'os conjectured that every graph of chromatic number at least $k$ contains a
subdivision of $K_k$ (a \emph{topological} $K_k$). The random graph refutes this
for large $k$, combining the clique-number concentration of Chapter~4 with a
Chernoff degree estimate.

\begin{theorem}{Erd\H{o}s--Fajtlowicz 1981}{hajos}
For all large $n$, the graph $G(n,1/2)$ has chromatic number
$\chi\sim n/(2\log_2 n)$ but its largest topological clique has order
$O(\sqrt n)$. In particular Haj\'os's conjecture fails.
\end{theorem}

\begin{proof}[Proof idea]
By Theorem~\ref{thm:clique-2pt} the independence number of $G(n,1/2)$ is
$\sim2\log_2 n$ \whp, so $\chi\ge n/\alpha\sim n/(2\log_2 n)$; a greedy argument
gives a matching upper bound. On the other hand a topological $K_t$ uses $t$
branch vertices pairwise joined by internally disjoint paths, so it needs
$\binom t2$ edge-disjoint connections among $t$ vertices --- and a Chernoff bound
on vertex degrees shows $G(n,1/2)$ simply lacks the connectivity to support
$t\gg\sqrt n$ branch vertices \whp. Since $n/\log n\gg\sqrt n$, the chromatic
number vastly exceeds the largest topological clique.
\end{proof}

\begin{insight}
The refutation splits across two chapters. The second moment fixes the
independence number, hence $\chi\sim n/(2\log_2 n)$; a Chernoff degree bound
caps the largest topological clique at $O(\sqrt n)$. Neither alone suffices, but
together they open a polynomial gap between the two quantities Haj\'os assumed
were linked. Sharp concentration of both parameters is exactly what makes
$G(n,1/2)$ a good place to hunt for counterexamples.
\end{insight}

% =====================================================================
\section*{Exercises}
\addcontentsline{toc}{section}{Exercises}
% =====================================================================

\begin{enumerate}[itemsep=4pt]
  \item Prove that with probability $1-o(1)$, every bipartite subgraph of
        $G(n,1/2)$ has at most $n^2/8+10n^{3/2}$ edges.
  \item (Unbalancing lights, upper bound.) Prove there is a constant $C$ so that
        for every $n$ there is a $\{-1,1\}$ matrix $A$ with
        $|y^\top A x|\le Cn^{3/2}$ for all $x,y\in\{-1,1\}^n$.
  \item Prove there is a constant $c>1$ so that for every $n$ there are at least
        $c^n$ points in $\mathbb R^n$ whose every triple spans a triangle with all
        angles below $61^\circ$.
  \item (Planted clique.) Give a deterministic polynomial-time algorithm that,
        on the union of $G(n,1/2)$ and an unlabelled clique of size
        $t=100\sqrt{n\log n}$, recovers a clique of size $t$ \whp.
  \item (Weighing coins.) You have $n$ coins of two known weights. A scale reports
        the total weight of any chosen subset. Show $O(n/\log n)$ weighings suffice
        to determine every coin's weight.
\end{enumerate}
